\begin{proof}[Proof of \cref{obj:thm:curve-upper}]
\textit{1. Constants and unit bound.} Write
\[
L:=\log(ed)=1+\log d,\qquad
\rho_{n,d}:=\frac{d}{nL},\qquad
\gamma_\epsilon:=3(1+\epsilon^{-1})+1.
\]
Let \(C_\epsilon^{\mathrm{proc}}>0\) be the constant supplied by \cref{obj:thm:integrated-process-certificate}, and set
\[
C_\epsilon:=\max\{1,C_\epsilon^{\mathrm{proc}},4096\gamma_\epsilon^2\}>0.
\]
Since \(d\ge2\), the inequality \(\log d\le d-1\) gives \(0<L\le d\), so \(\rho_{n,d}\ge0\).

The interval \(I=[b_0,1-b_0]\) is nonempty and contained in \([0,1]\). Write
\[
p_j=P(X=j),\qquad
\mu_{aj}=\begin{cases}E_P[Y(a)\mid X=j],&p_j>0,\\0,&p_j=0,\end{cases}
\qquad
\tau_j=\mu_{1j}-\mu_{0j},\qquad
B(P)=\sum_{j=1}^d p_j\mu_{0j}.
\]
The binary potential-outcome condition in \cref{obj:def:model-class} gives \(0\le\mu_{aj}\le1\), including the stated convention at empty cells. For any \(b\in I\), the zero policy is feasible under \cref{obj:def:budget-policy-class}, and the value of every feasible policy is
\[
B(P)+\sum_{j=1}^d p_j\pi_j\tau_j
=\sum_{j=1}^d p_j\bigl\{(1-\pi_j)\mu_{0j}+\pi_j\mu_{1j}\bigr\}\in[0,1].
\]
Hence \(V_b(P)\in[0,1]\) by \cref{obj:def:budget-value}. The estimator in \cref{obj:def:jf-estimator} is either \(1/2\) or clipped to \([0,1]\). Under \cref{obj:ass:iid-sampling}, its squared supremum loss therefore satisfies
\[
E_P\!\left[\sup_{b\in I}
 \bigl|\widehat V_b^{\mathrm{JF}}-V_b(P)\bigr|^2\right]\le1.
\]
% lean: CausalSmith.Stat.DiscreteBudgetvalueCurve.jfEstimator_curveRisk_le_one

\textit{2. Alphabets with \(2\le d<16\).} In this branch, \(\widehat F\) is the empirical threshold process of \cref{obj:def:jf-estimator}. It is bounded on \([0,1]\) by \cref{obj:lem:empirical-threshold-process-bounded}. The dual representation and its uniform perturbation bound in \cref{obj:prop:dual-representation}, together with contraction by clipping toward \(V_b(P)\in[0,1]\), give, sample by sample,
\[
\sup_{b\in I}
 \bigl|\widehat V_b^{\mathrm{JF}}-V_b(P)\bigr|^2
\le
\left(\sup_{\lambda\in[0,1]}
 \bigl|\widehat F(\lambda)-F_P(\lambda)\bigr|\right)^2.
\]
% lean: CausalSmith.Stat.DiscreteBudgetvalueCurve.jfEstimator_curveLoss_le_processError

Here is the pointwise estimate used for the empirical process. Let \(\mathcal Z=\{0,1\}^2\), index \(\zeta=(a,y)\), and write
\[
q_{\zeta j}=P(X=j,A=a,Y=y),\qquad
s_a(u)=u_{(a,0)}+u_{(a,1)},\qquad
s(u)=s_0(u)+s_1(u)
\]
for every nonnegative four-vector \(u=(u_\zeta)_{\zeta\in\mathcal Z}\). For nonnegative four-vectors \(u,v\), define
\[
\Delta(u,v):=\sum_{\zeta\in\mathcal Z}|u_\zeta-v_\zeta|.
\]
For the formula in \cref{obj:def:dual-process}, explicitly define its value at the zero vector by \(g_a(0)=0\). This is its continuous extension, since \(0\le g_a(u)\le s(u)\) for every nonnegative \(u\) with \(s(u)>0\). Thus \(f_\lambda(0)=0\), and the stabilized arm contributions, including those of empty covariate cells, obey
\[
|g_a(u)-g_a(v)|\le(1+\epsilon^{-1})\Delta(u,v),
\qquad
|s(u)-s(v)|\le\Delta(u,v).
\]
To see the first inequality, an arm with \(s_a(u)\le\epsilon s(u)\) has \(g_a(u)=u_{a1}/\epsilon\); an arm with \(s_a(u)\ge\epsilon s(u)>0\) has \(g_a(u)=s(u)u_{a1}/s_a(u)\). When both vectors are in the latter regime, use
\[
\left|\frac{u_{a1}}{s_a(u)}-\frac{v_{a1}}{s_a(v)}\right|
\le
\frac{|u_{a1}-v_{a1}|+|u_{a0}-v_{a0}|}{s_a(u)},
\qquad
\frac{s(u)}{s_a(u)}\le\epsilon^{-1},
\]
and \(u_{a1}/s_a(u),v_{a1}/s_a(v)\in[0,1]\). When both vectors are in the first regime, \(|g_a(u)-g_a(v)|=|u_{a1}-v_{a1}|/\epsilon\le\epsilon^{-1}\Delta(u,v)\). When \(u\) is in the first regime and \(v\) in the second, use
\[
g_a(u)-g_a(v)
=
\frac{u_{a1}-v_{a1}}{\epsilon}
+
\frac{v_{a1}}{s_a(v)}
 \frac{s_a(v)-\epsilon s(v)}{\epsilon}.
\]
The second term is nonnegative and \(0\le v_{a1}/s_a(v)\le1\). Since \(s_a(u)\le\epsilon s(u)\), its numerator satisfies
\[
s_a(v)-\epsilon s(v)
\le s_a(v)-s_a(u)+\epsilon\{s(u)-s(v)\}.
\]
Combining these inequalities cancels the \(a1\) coordinate in the upper bound and gives
\[
-\epsilon^{-1}\Delta(u,v)
\le g_a(u)-g_a(v)
\le \epsilon^{-1}(v_{a0}-u_{a0})+s(u)-s(v)
\le(1+\epsilon^{-1})\Delta(u,v).
\]
Swapping \(u\) and \(v\) handles the reverse case. If either total mass is zero, that vector is zero and \(0\le g_a(w)\le s(w)\) gives the bound directly. Since \(x\mapsto[x]_+\) is one-Lipschitz and \(0\le\lambda\le1\), it follows that
\[
|f_\lambda(u)-f_\lambda(v)|
\le
2|g_0(u)-g_0(v)|+|g_1(u)-g_1(v)|
 +\lambda|s(u)-s(v)|
\le\gamma_\epsilon\Delta(u,v).
\]
% lean: CausalSmith.Stat.DiscreteBudgetvalueCurve.budget_thresholdFunReal_pointwise_lipschitz

Applying this cellwise and using Cauchy--Schwarz across the \(4d\) category-cell coordinates yields
\[
\left(\sup_{\lambda\in[0,1]}
 |\widehat F(\lambda)-F_P(\lambda)|\right)^2
\le
4d\,\gamma_\epsilon^2
\sum_{j=1}^d\sum_{\zeta\in\mathcal Z}
 (\widehat q_{\zeta j}-q_{\zeta j})^2.
\]
Each \(\widehat q_{\zeta j}\) is the empirical mass of one observed atom. Its Bernoulli second moment under \cref{obj:ass:iid-sampling} gives
\[
E_P(\widehat q_{\zeta j}-q_{\zeta j})^2
=\frac{q_{\zeta j}(1-q_{\zeta j})}{n}\le\frac1n.
\]
Consequently,
\[
E_P\!\left[\sup_{b\in I}
 |\widehat V_b^{\mathrm{JF}}-V_b(P)|^2\right]
\le
\frac{16d^2\gamma_\epsilon^2}{n}.
\]
% lean: CausalSmith.Stat.DiscreteBudgetvalueCurve.jfEstimator_empiricalCurveRisk_le

Because \(L\le d<16\), we have \(dL\le256\), and hence
\[
\frac{16d^2\gamma_\epsilon^2}{n}
\le4096\gamma_\epsilon^2\rho_{n,d}.
\]
If \(\rho_{n,d}\ge1\), the unit bound above is at most \(C_\epsilon\min\{1,\rho_{n,d}\}\). If \(\rho_{n,d}<1\), the displayed empirical bound is at most the same quantity.
% lean: CausalSmith.Stat.DiscreteBudgetvalueCurve.empiricalCurveRate_algebra

\textit{3. Alphabets with \(d\ge16\) and \(d/L\le n\).} Here \(\rho_{n,d}\le1\). The estimator uses the conditionally averaged threshold process \(\widehat F\) specified in \cref{obj:def:jf-estimator}. It is bounded on \([0,1]\) by \cref{obj:lem:averaged-threshold-process-bounded}. The same dual perturbation and clipping argument bounds its curve risk by the averaged process risk:
\[
E_P\!\left[\sup_{b\in I}
 |\widehat V_b^{\mathrm{JF}}-V_b(P)|^2\right]
\le
E_P\!\left[
 \left(\sup_{\lambda\in[0,1]}
 |\widehat F(\lambda)-F_P(\lambda)|\right)^2\right].
\]
% lean: CausalSmith.Stat.DiscreteBudgetvalueCurve.jfEstimator_averagedCurveRisk_le

Under the stated sampling and model conditions, \cref{obj:thm:integrated-process-certificate} bounds the right-hand side by \(C_\epsilon^{\mathrm{proc}}\rho_{n,d}\). Thus the curve risk is at most
\[
C_\epsilon^{\mathrm{proc}}\rho_{n,d}
\le C_\epsilon\min\{1,\rho_{n,d}\}.
\]
% lean: CausalSmith.Stat.DiscreteBudgetvalueCurve.integrated_process_certificate

\textit{4. Alphabets with \(d\ge16\) and \(n<d/L\).} Now \(\rho_{n,d}>1\), and \cref{obj:def:jf-estimator} sets the entire estimated curve equal to \(1/2\). The unit bound above gives
\[
E_P\!\left[\sup_{b\in I}
 |\widehat V_b^{\mathrm{JF}}-V_b(P)|^2\right]
\le1\le C_\epsilon\min\{1,\rho_{n,d}\}.
\]
% lean: CausalSmith.Stat.DiscreteBudgetvalueCurve.jfEstimator_curveRisk_le_one

These cases give the claimed inequality for every \(P\in\mathcal M_{d,\epsilon}\), with one constant depending only on \(\epsilon\). Consider the set of expected squared losses indexed by these laws. If this set is nonempty, its supremum obeys the same bound because every member does. If it is empty, define its supremum as zero; the bound then follows from \(C_\epsilon>0\) and \(\rho_{n,d}\ge0\). In the present model the set is nonempty: take \(X\) uniform on \([d]\), take \(A\) independent of \(X\) with treatment probability \(1/2\), and set \(Y(0)=Y(1)=Y=0\) almost surely. This law satisfies consistency and conditional exchangeability, and its treatment propensity in every cell is \(1/2\), as required by \(0<\epsilon<1/2\). Thus the nonempty case gives the stated supremum bound.
% lean: CausalSmith.Stat.DiscreteBudgetvalueCurve.curve_upper
\end{proof}
